\documentclass[c]{beamer}

\usepackage[utf8]{inputenc}
\usepackage[french]{babel}

\usepackage{amsfonts}
\usepackage{amsmath}

\newtheorem*{deffr}{Définition}
\newtheorem*{propriete}{Propriété}
\newtheorem*{theofr}{Théorème}

\usepackage{graphicx}
\usepackage{array}

\usetheme{Warsaw}

\title{Dynamique d'un graphe}
\author{Igor Colin}
\date{\today}

\setbeamertemplate{navigation symbols}{}

\AtBeginSection[]
{
    \begin{frame}<beamer>
        \tableofcontents[currentsection]
    \end{frame}
}

\begin{document}

%\maketitle

\section{Objectifs}

\begin{frame}
    \frametitle{Objectifs}

    \begin{itemize}
        \item<1-> Modélisation des données à l'aide d'un graphe
            \begin{itemize}
                \item N\oe{}uds : utilisateurs
                \item Arrêtes : contact, influence, etc.
            \end{itemize}
        \item<2-> Intérêt pour la notion de communauté :
            \begin{itemize}
                \item Suggestion de contenus adaptés
                \item Proposition de groupe
            \end{itemize}
        \item<3-> Calculs sur serveur mais aussi en local
            \begin{itemize}
                \item Puissance de calcul raisonnable
                \item Mémoire limitée
            \end{itemize}
    \end{itemize}
\end{frame}

\section{Génération aléatoire de graphe}

\begin{frame}
    \frametitle{Motivation}

    \begin{itemize}
        \item<1-> Volume de données faible
        \item<2-> Méthodes robustes
        \item<3-> Tests de scénarios
    \end{itemize}
\end{frame}

\begin{frame}
    \frametitle{Modèle $p_1$}
    \framesubtitle{Idée}

    \begin{itemize}
        \item<1-> Principe
            \begin{itemize}
                \item Modèle paramétrique
                \item Paramètres physiques/concrets (popularité, etc.)
                \item Apprentissage des paramètres sur un réseau réel (maximum
                    de vraisemblance)
            \end{itemize}
        \item<2-> Paramètres
            \begin{itemize}
                \item Expansion global du graphe $\theta$
                \item Expansion d'un n\oe{}ud $\alpha_i$
                \item Popularité d'un n\oe{}ud $\beta_i$
                \item Intérêt mutuel $\rho_{ij} = \rho_{ji}$ (ex: Twitter)
            \end{itemize}
    \end{itemize}
\end{frame}

\begin{frame}
    \frametitle{Modèle $p_1$}
    \framesubtitle{Vraisemblance}

    \begin{table}
        \begin{tabular}[c]{m{.2\textwidth} m{.05\textwidth} m{.5\textwidth}}
                \includegraphics[width=.2\textwidth]{./figures/p1_00.pdf}
                & : &
                $\lambda_{ij}$ (normalisation)\\
                \includegraphics[width=.2\textwidth]{./figures/p1_10.pdf}
                & : &
                $\lambda_{ij} + \alpha_i + \beta_j + \theta$ \\
                \includegraphics[width=.2\textwidth]{./figures/p1_01.pdf}
                & : &
                $\lambda_{ij} + \alpha_j + \beta_i + \theta$ \\
                \includegraphics[width=.2\textwidth]{./figures/p1_11.pdf}
                & : &
                $\lambda_{ij} + \alpha_i + \alpha_j + \beta_i + \beta_j + 2\theta + \rho_{ij}$
        \end{tabular}
        \caption{Log-vraisemblance en fonction des paramètres.}
    \end{table}

\end{frame}

\begin{frame}
    \frametitle{Modèle $p_1$}
    \framesubtitle{Exemple}
    \begin{figure}
        \includegraphics[width=.7\textwidth]{./figures/p1_popularite.png}
        \caption{Modifications du graphe pour différentes valeurs de $\beta_0$.}
    \end{figure}
\end{frame}

\begin{frame}
    \frametitle{Modèle $p_1$}
    \framesubtitle{Discussion}

    Avantages du modèle
    \begin{itemize}
        \item Modèle simple
        \item Calcul aisé dans cas favorables
        \item Paramètres interprétables
        \item Nombreuses variantes :
            \begin{itemize}
                \item Dynamique (chaîne de Markov)
                \item Bayésienne
            \end{itemize}
    \end{itemize}
\end{frame}

\begin{frame}
    \frametitle{Modèle $p_1$}
    \framesubtitle{Discussion}

    \begin{columns}
        \begin{column}{.6\textwidth}
             \'Evaluation du modèle : quelle métrique utiliser ?
            \begin{itemize}
                \item Métriques usuelles peu pertinentes (précision)
                \item Pertinence de la distance directement liée à l'exploitation
                    souhaitée des données
            \end{itemize}
        \end{column}
        \begin{column}{.5\textwidth}
            \begin{figure}
                \includegraphics[width=.9\textwidth]{./figures/p1_distance.png}
            \end{figure}
        \end{column}
    \end{columns}

\end{frame}

%\section{Détection de communauté}
%
%\begin{frame}
%    \frametitle{Approche spectrale usuelle}
%
%    \begin{itemize}
%        \item Objectif : trouver le partionnement du graphe qui maximise la
%            modularité
%        \item Obtention via la matrice laplacienne
%        \item Calcul de valeurs propres coûteux : obligation d'effectuer le
%            calcul sur un serveur ou de revoir complétement la méthode si le
%            réseau est trop grand
%    \end{itemize}
%\end{frame}
%
%\begin{frame}
%    \frametitle{Approche spectrale randomisée}
%
%
%\end{frame}

\end{document}
